% META: {"id":"TSPE-GEOESPACE-ME-001","chapitre":"TSPE-GEOMETRIE-ESPACE","type_objet":"methode","capacites_codes":["C13","C16"],"methodes":["M1"],"sources_inspiration":[],"status":"approved","origin":"generated","status_history":["generated","audited","accepted","approved"],"release_acceptance":"RELEASE_OWNER_BATCH_ACCEPTANCE","acceptance_closure_digest":"sha256:9b3ccf9a81c5520fbb7e03b7b2d3e7bf2057a02834b6d908bf3be5f49f3b553f","acceptance_receipt_digest":"sha256:4fad86f0282e0fa4e0419cca1ad6379ae7af5a280c04a4a90880f90a1c044242"}
% BEGIN-VERIFY
% from sympy import Matrix
% A=Matrix([2,-1,3]); n=Matrix([1,4,-2])
% d = -(n.dot(A))
% assert d == 8
% assert n.dot(A) + d == 0          # A appartient au plan x+4y-2z+8=0
% assert n.dot(Matrix([0,0,0])) + d != 0   # l'origine n'y est pas
% END-VERIFY
\begin{fichemethode}{M1}{Déterminer l'équation cartésienne d'un plan}

\quandUtiliser{%
  Pour écrire l'équation cartésienne d'un plan connaissant un point et un
  vecteur normal, pour démontrer cette équation, et inversement pour lire un
  vecteur normal à partir d'une équation donnée.
}

\pasApas{%
  \begin{enumerate}
    \item \textbf{Identifier} un point $A(x_A;y_A;z_A)$ du plan et un vecteur
          normal $\vec{n}(a;b;c)$. Si $\vec{n}$ n'est pas donné, le déterminer
          en écrivant qu'il est orthogonal à deux vecteurs non colinéaires du
          plan, puis en résolvant le système obtenu.
    \item \textbf{Écrire} la condition d'appartenance
          $\vec{AM}\cdot\vec{n}=0$, soit
          $a(x-x_A)+b(y-y_A)+c(z-z_A)=0$.
    \item \textbf{Développer} pour obtenir la forme réduite
          $ax+by+cz+d=0$, avec $d=-(a x_A+b y_A+c z_A)$.
    \item \textbf{Lire un vecteur normal} à partir d'une équation
          $ax+by+cz+d=0$~: c'est directement $\vec{n}(a;b;c)$.
  \end{enumerate}
}

\verifier{%
  Substituer les coordonnées de $A$ dans l'équation obtenue~: elle doit donner
  $0$. Si un second point du plan est connu, le contrôler également. Vérifier
  enfin qu'un vecteur du plan est bien orthogonal à $\vec{n}$.
}

\pieges{%
  Ne pas oublier le terme constant~: $ax+by+cz=0$ est l'équation du plan
  passant par l'\emph{origine}, parallèle au plan cherché mais distinct de lui.
  Attention aussi au signe de $d$, qui vient de la distribution du signe moins.
  Enfin, $\vec{n}(a;b;c)$ est \emph{normal} au plan~: il ne le dirige pas.
}

\exempleRedige{%
  Soit $A(2;-1;3)$ et $\vec{n}(1;4;-2)$. La condition
  $\vec{AM}\cdot\vec{n}=0$ s'écrit
  $1(x-2)+4(y+1)-2(z-3)=0$, soit $x-2+4y+4-2z+6=0$ et enfin
  \[
    x+4y-2z+8=0 .
  \]
  Contrôle~: $2+4\times(-1)-2\times3+8=2-4-6+8=0$, donc $A$ appartient bien au
  plan.
}

\end{fichemethode}
